\chapter{Chapter 3: Partial Derivatives}

\begin{enumerate}

%%% FUNCTIONS OF SEVERAL VARIABLES (Problems 1--10) %%%

% Problem 1
\item \textbf{Solution.}
We have $M = f(P,r,t) = P\dfrac{\frac{r}{12}\left(1+\frac{r}{12}\right)^{12t}}{\left(1+\frac{r}{12}\right)^{12t}-1}$.

Substituting $P = 300{,}000$, $r = 0.06$, $t = 30$:
\[
\frac{r}{12} = \frac{0.06}{12} = 0.005, \qquad 12t = 360.
\]
\[
\left(1.005\right)^{360} \approx 6.02258.
\]
Therefore:
\[
M = 300{,}000 \cdot \frac{0.005 \times 6.02258}{6.02258 - 1} = 300{,}000 \cdot \frac{0.030113}{5.02258} \approx 300{,}000 \times 0.005996 \approx \boxed{\$1{,}798.65}.
\]

% Problem 2
\item \textbf{Solution.}
Setting $f(x,y) = x^2 + y^2 = k$:
\begin{itemize}
    \item $k = 1$: $x^2 + y^2 = 1$ is a circle of radius 1 centered at the origin.
    \item $k = 2$: $x^2 + y^2 = 2$ is a circle of radius $\sqrt{2}$ centered at the origin.
    \item $k = 4$: $x^2 + y^2 = 4$ is a circle of radius 2 centered at the origin.
\end{itemize}
All level curves are concentric circles centered at the origin. The sketch shows three circles of radii 1, $\sqrt{2}$, and 2.

% Problem 3
\item \textbf{Solution.}
The natural logarithm requires its argument to be strictly positive:
\[
x^2 + y^2 - 4 > 0 \implies x^2 + y^2 > 4.
\]
The domain is the set of all points $(x,y)$ \emph{outside} the circle of radius 2 centered at the origin:
\[
\boxed{D = \{(x,y) : x^2 + y^2 > 4\}}.
\]

% Problem 4
\item \textbf{Solution.}
Setting $f(x,y,z) = x^2 + y^2 + z^2 = k$ for $k \geq 0$:
\begin{itemize}
    \item If $k = 0$, the level ``surface'' is the single point $(0,0,0)$.
    \item If $k > 0$, the level surface is a sphere of radius $\sqrt{k}$ centered at the origin.
\end{itemize}
The level surfaces are concentric spheres centered at the origin.

% Problem 5
\item \textbf{Solution.}
For the square root to be defined, we need $9 - x^2 - y^2 \geq 0$, i.e., $x^2 + y^2 \leq 9$.

\textbf{Domain:} The closed disk of radius 3 centered at the origin:
$D = \{(x,y) : x^2 + y^2 \leq 9\}$.

\textbf{Range:} The minimum value of $f$ is 0 (when $x^2 + y^2 = 9$) and the maximum value is 3 (when $x^2 + y^2 = 0$). Thus the range is $\boxed{[0,3]}$.

% Problem 6
\item \textbf{Solution.}
The parametric curve gives $x = t$, $y = t^2$, $z = t^3$. On the surface $z = y^{3/2}$ we need:
\[
t^3 = (t^2)^{3/2} = |t|^3.
\]
\begin{itemize}
    \item If $t \geq 0$: $t^3 = t^3$, which holds for all $t \geq 0$.
    \item If $t < 0$: $t^3 = -|t|^3$ while $|t|^3 > 0$, so $t^3 \neq |t|^3$.
\end{itemize}
Therefore the curve lies on the surface for $\boxed{t \geq 0}$.

% Problem 7
\item \textbf{Solution.}
We need $\dfrac{x}{y-1} \geq 0$. This requires $x$ and $y-1$ to have the same sign (or $x = 0$):
\begin{itemize}
    \item $x \geq 0$ and $y > 1$ (both non-negative, denominator nonzero), or
    \item $x \leq 0$ and $y < 1$ (both non-positive, denominator nonzero).
\end{itemize}
Note $y \neq 1$ since that makes the denominator zero.

The domain is $D = \{(x,y) : x \geq 0,\, y > 1\} \cup \{(x,y) : x \leq 0,\, y < 1\}$.

The sketch shows the first and third quadrants relative to the point $(0,1)$: the region in the upper half-plane ($y > 1$) for $x \geq 0$, and the lower half-plane ($y < 1$) for $x \leq 0$, excluding the line $y = 1$.

% Problem 8
\item \textbf{Solution.}
Answers will vary. Suppose the room is rectangular with length $L$, width $W$, and height $H$ (all in meters). For example, $L = 5$, $W = 4$, $H = 3$.

The total wall area is $A = 2H(L + W)$.

If the paint costs $\$P$ per square meter, and we apply one coat of primer plus one coat of paint (both at the same coverage rate), the cost function is:
\[
C(P, L, W, H) = 2P \cdot 2H(L + W) = 4H(L+W)P.
\]
If primer costs a different amount $\$P_{\text{primer}}$ per square meter, then:
\[
C = 2H(L+W)(P + P_{\text{primer}}).
\]
Using the example dimensions: $A = 2(3)(5+4) = 54$ m$^2$, and the cost is $54(P + P_{\text{primer}})$ dollars.

% Problem 9
\item \textbf{Solution.}
Given $T = \dot{m}v_e + (P_e - P_a)A_e$ with $P_e = 101$ kPa, $P_a = 100$ kPa, $A_e = 2$ m$^2$:
\[
(P_e - P_a)A_e = (101 - 100)(1000)(2) = 2000 \text{ N} = 2 \text{ kN}.
\]
Solving for $\dot{m}$:
\[
T = \dot{m}v_e + 2000 \implies \boxed{\dot{m} = \frac{T - 2000}{v_e}}.
\]

% Problem 10
\item \textbf{Solution.}
The profit function simplifies to:
\[
P = f(x,y) = (4x + 5y) - (2x + 3y) - 100 = 2x + 2y - 100.
\]
\[
f(50,30) = 2(50) + 2(30) - 100 = 100 + 60 - 100 = \boxed{\$60}.
\]

With 6\% taxes applied to revenue (taxes apply to income, not to debits/costs):
\[
P_{\text{tax}} = (4x + 5y)(1 - 0.06) - (2x + 3y) - 100 = 0.94(4x+5y) - (2x+3y) - 100
\]
\[
= 3.76x + 4.70y - 2x - 3y - 100 = 1.76x + 1.70y - 100.
\]

%%% LIMITS AND CONTINUITY (Problems 11--19) %%%

% Problem 11
\item \textbf{Solution.}
Test along $y = 0$: $\displaystyle\lim_{x \to 0} \frac{x \cdot 0}{x^2 + 0} = 0$.

Test along $y = x$: $\displaystyle\lim_{x \to 0} \frac{x \cdot x}{x^2 + x^2} = \lim_{x \to 0} \frac{x^2}{2x^2} = \frac{1}{2}$.

Since the limits along different paths are different ($0 \neq \frac{1}{2}$), the limit \textbf{does not exist}.

% Problem 12
\item \textbf{Solution.}
Convert to polar coordinates: $x = r\cos\theta$, $y = r\sin\theta$:
\[
\frac{x^2 y}{x^2 + y^2} = \frac{r^2\cos^2\theta \cdot r\sin\theta}{r^2} = r\cos^2\theta\sin\theta.
\]
Then:
\[
|r\cos^2\theta\sin\theta| \leq r \cdot 1 \cdot 1 = r \to 0 \text{ as } r \to 0.
\]
By the Squeeze Theorem, $\boxed{\displaystyle\lim_{(x,y)\to(0,0)} \frac{x^2y}{x^2+y^2} = 0}$.

% Problem 13
\item \textbf{Solution.}
Test along $y = 0$: $\displaystyle\lim_{x \to 0} \frac{x \cdot 0}{x^2 + 0} = 0$.

Test along $x = y^2$:
\[
\lim_{y \to 0} \frac{y^2 \cdot y^2}{(y^2)^2 + y^4} = \lim_{y \to 0} \frac{y^4}{2y^4} = \frac{1}{2}.
\]

Since the limits along different paths differ ($0 \neq \frac{1}{2}$), $\displaystyle\lim_{(x,y)\to(0,0)} f(x,y)$ does not exist. Therefore $f$ is \textbf{not continuous} at $(0,0)$.

% Problem 14
\item \textbf{Solution.}
Along $y = 0$: $\displaystyle\lim_{x \to 0} \frac{x^2 - 0}{x^2 + 0} = 1$.

Along $x = 0$: $\displaystyle\lim_{y \to 0} \frac{0 - y^2}{0 + y^2} = -1$.

Since $1 \neq -1$, the limit \textbf{does not exist}.

% Problem 15
\item \textbf{Solution.}
We show $\displaystyle\lim_{(x,y)\to(0,0)} \frac{x^3 y}{x^4 + y^2} = 0$.

We use the AM--GM inequality: $x^4 + y^2 \geq 2x^2|y|$ (since $a + b \geq 2\sqrt{ab}$ for $a,b \geq 0$).

Therefore:
\[
\left|\frac{x^3 y}{x^4 + y^2}\right| \leq \frac{|x|^3|y|}{2x^2|y|} = \frac{|x|}{2} \to 0 \text{ as } (x,y) \to (0,0).
\]
(When $y = 0$, the expression equals 0 directly.)

By the Squeeze Theorem, $\displaystyle\lim_{(x,y)\to(0,0)} f(x,y) = 0 = f(0,0)$, so $f$ is continuous at $(0,0)$. \qed

% Problem 16
\item \textbf{Solution.}
Let $u = xy - 2$. As $(x,y) \to (1,2)$, $u \to 0$. Write:
\[
\frac{\sin(xy - 2)}{x - 1}.
\]
Substituting $x = 1 + h$ (so $h \to 0$):
\[
\frac{\sin((1+h)y - 2)}{h} = \frac{\sin(y - 2 + hy)}{h}.
\]
As $h \to 0$ and $y \to 2$: the argument $y - 2 + hy \to 0 + 0 = 0$. Using the approximation $\sin(\alpha) \approx \alpha$ for small $\alpha$:
\[
\frac{\sin(y - 2 + hy)}{h} \approx \frac{y - 2 + hy}{h}.
\]
Along the path $y = 2$ (taking the limit with $y = 2$):
\[
\frac{\sin(2h)}{h} \to 2.
\]
More rigorously, set $y = 2$ first: $\displaystyle\lim_{x \to 1} \frac{\sin(2x - 2)}{x-1} = \lim_{h \to 0} \frac{\sin(2h)}{h} = 2$.

For a general path, write $y = 2 + k$ where $k \to 0$:
\[
\frac{\sin(k + (1+h)(2+k) - 2)}{h} = \frac{\sin(k + 2h + hk)}{h}.
\]
Note that $(1+h)(2+k) - 2 = 2 + k + 2h + hk - 2 = k + 2h + hk$. Therefore:
\[
\frac{\sin(k + 2h + hk)}{h} = \frac{\sin(k + 2h + hk)}{k + 2h + hk} \cdot \frac{k + 2h + hk}{h}.
\]
The first factor $\to 1$. The second factor is $\dfrac{k}{h} + 2 + k$. For the limit to exist, we need $k/h$ to have a definite limit; but if $k = 0$ (i.e., $y \to 2$ at the same rate), this gives 2. Along $k = h$: $(h/h) + 2 + h \to 3$. The limits differ, so the two-variable limit \textbf{does not exist} in general.

However, if the problem intends $y \to 2$ independently (as a standard multivariable limit), the limit does not exist. If $y = 2$ is fixed and only $x \to 1$, then the answer is $\boxed{2}$.

% Problem 17
\item \textbf{Solution.}
We have $\phi(x,y) = \dfrac{Q}{2\pi}\ln\sqrt{x^2+y^2} = \dfrac{Q}{4\pi}\ln(x^2+y^2)$.

Let $r = \sqrt{x^2 + y^2}$. As $(x,y) \to (0,0)$, $r \to 0^+$, so:
\[
\ln(x^2 + y^2) = \ln(r^2) = 2\ln r \to -\infty.
\]
Therefore (assuming $Q > 0$):
\[
\phi(x,y) = \frac{Q}{4\pi}\ln(r^2) = \frac{Q}{2\pi}\ln r \to -\infty.
\]
This holds regardless of the path, since $r \to 0$ forces $\ln r \to -\infty$. Thus $\displaystyle\lim_{(x,y)\to(0,0)} \phi(x,y) = -\infty$. \qed

% Problem 18
\item \textbf{Solution.}
The magnitude of the electric field at $(x,y)$ with source at $(x_0, y_0)$ is:
\[
|\mathbf{E}(x,y)| = \frac{|q|}{4\pi\epsilon_0} \cdot \frac{\sqrt{(x-x_0)^2 + (y-y_0)^2}}{[(x-x_0)^2+(y-y_0)^2]^{3/2}} = \frac{|q|}{4\pi\epsilon_0} \cdot \frac{1}{(x-x_0)^2+(y-y_0)^2}.
\]
Let $r = \sqrt{(x-x_0)^2+(y-y_0)^2}$, so $|\mathbf{E}| = \dfrac{|q|}{4\pi\epsilon_0 r^2}$.

As $(x,y) \to (x_0, y_0)$: $r \to 0$ and so $|\mathbf{E}| \to +\infty$.

Taking the limit as $(x,y) \to (0,0)$ with the charge at $(x_0, y_0)$:

If $(x_0, y_0) \neq (0,0)$, then $r \to \sqrt{x_0^2 + y_0^2} > 0$ and
\[
\lim_{(x,y)\to(0,0)} |\mathbf{E}(x,y)| = \frac{|q|}{4\pi\epsilon_0(x_0^2+y_0^2)}.
\]

If $(x_0, y_0) = (0,0)$, then $\displaystyle\lim_{(x,y)\to(0,0)} |\mathbf{E}(x,y)| = +\infty$.

% Problem 19
\item \textbf{Solution.}
We have $\sigma(r,\theta) = \dfrac{A}{r^2} + Br^2$ where $r = \sqrt{x^2+y^2}$.

As $(x,y) \to (0,0)$, $r \to 0^+$:
\begin{itemize}
    \item If $A \neq 0$: $\dfrac{A}{r^2} \to \pm\infty$ (depending on the sign of $A$), and $Br^2 \to 0$.
\end{itemize}
Therefore, if $A \neq 0$, the limit diverges: $\sigma \to +\infty$ if $A > 0$ and $\sigma \to -\infty$ if $A < 0$.

If $A = 0$: $\sigma = Br^2 \to 0$.

In general (for a physically meaningful stress function with $A \neq 0$), the limit \textbf{does not exist} as a finite value; $\sigma \to \pm\infty$.

%%% PARTIAL DERIVATIVES (Problems 20--30) %%%

% Problem 20
\item \textbf{Solution.}
$f(x,y) = x^2y + 3xy^2$.
\[
f_x = \frac{\partial}{\partial x}(x^2y + 3xy^2) = 2xy + 3y^2.
\]
\[
f_y = \frac{\partial}{\partial y}(x^2y + 3xy^2) = x^2 + 6xy.
\]

% Problem 21
\item \textbf{Solution.}
$f(x,y) = \sin(xy)e^x$.

\textbf{First partials:}
\[
f_x = y\cos(xy)e^x + \sin(xy)e^x = e^x[y\cos(xy) + \sin(xy)].
\]
\[
f_y = x\cos(xy)e^x.
\]

\textbf{Second partials:}
\begin{align*}
f_{xx} &= \frac{\partial}{\partial x}\bigl[e^x(y\cos(xy)+\sin(xy))\bigr] \\
&= e^x[y\cos(xy)+\sin(xy)] + e^x[-y^2\sin(xy)+y\cos(xy)] \\
&= e^x[2y\cos(xy) + \sin(xy) - y^2\sin(xy)].
\end{align*}

\begin{align*}
f_{yy} &= \frac{\partial}{\partial y}[x\cos(xy)e^x] = -x^2\sin(xy)e^x.
\end{align*}

\begin{align*}
f_{xy} &= \frac{\partial}{\partial y}\bigl[e^x(y\cos(xy)+\sin(xy))\bigr] \\
&= e^x[\cos(xy) - xy\sin(xy) + x\cos(xy)] \\
&= e^x[(1+x)\cos(xy) - xy\sin(xy)].
\end{align*}

\begin{align*}
f_{yx} &= \frac{\partial}{\partial x}[x\cos(xy)e^x] \\
&= \cos(xy)e^x + x(-y\sin(xy))e^x + x\cos(xy)e^x \\
&= e^x[(1+x)\cos(xy) - xy\sin(xy)].
\end{align*}

We verify $f_{xy} = f_{yx}$, consistent with Clairaut's theorem.

% Problem 22
\item \textbf{Solution.}
$f(x,y,z) = xyz + e^{xy}$.

\textbf{First partials:}
\[
f_x = yz + ye^{xy}, \quad f_y = xz + xe^{xy}, \quad f_z = xy.
\]

\textbf{Mixed second partials:}
\[
f_{xy} = z + e^{xy} + xye^{xy}, \qquad f_{yx} = z + e^{xy} + xye^{xy}.
\]
\[
f_{xz} = y, \qquad f_{zx} = y.
\]
\[
f_{yz} = x, \qquad f_{zy} = x.
\]

All mixed partials satisfy $f_{xy} = f_{yx}$, $f_{xz} = f_{zx}$, $f_{yz} = f_{zy}$.

% Problem 23
\item \textbf{Solution.}
$f(x,y) = x^2e^y + \ln(xy) = x^2e^y + \ln x + \ln y$.
\[
f_x = 2xe^y + \frac{1}{x}, \qquad f_y = x^2e^y + \frac{1}{y}.
\]
\[
f_{xy} = \frac{\partial}{\partial y}\left(2xe^y + \frac{1}{x}\right) = 2xe^y.
\]
\[
f_{yx} = \frac{\partial}{\partial x}\left(x^2e^y + \frac{1}{y}\right) = 2xe^y.
\]
Since $f_{xy} = f_{yx} = 2xe^y$, Clairaut's theorem is verified. \qed

% Problem 24
\item \textbf{Solution.}
$f(x,y) = \arctan\!\left(\dfrac{y}{x}\right)$.
\[
f_x = \frac{1}{1+(y/x)^2} \cdot \left(-\frac{y}{x^2}\right) = \frac{-y}{x^2+y^2}.
\]
\[
f_y = \frac{1}{1+(y/x)^2} \cdot \frac{1}{x} = \frac{x}{x^2+y^2}.
\]

\textbf{Second partials:}
\[
f_{xx} = \frac{\partial}{\partial x}\left(\frac{-y}{x^2+y^2}\right) = \frac{-y \cdot (-2x)}{(x^2+y^2)^2} \cdot (-1) \cdot (-1) = \frac{2xy}{(x^2+y^2)^2}.
\]
Let us compute carefully: $f_x = -y(x^2+y^2)^{-1}$, so
\[
f_{xx} = -y \cdot (-1)(x^2+y^2)^{-2} \cdot 2x = \frac{2xy}{(x^2+y^2)^2}.
\]
\[
f_{yy} = \frac{\partial}{\partial y}\left(\frac{x}{x^2+y^2}\right) = x \cdot (-1)(x^2+y^2)^{-2}\cdot 2y = \frac{-2xy}{(x^2+y^2)^2}.
\]
\[
f_{xy} = \frac{\partial}{\partial y}\left(\frac{-y}{x^2+y^2}\right) = \frac{-(x^2+y^2) + y \cdot 2y}{(x^2+y^2)^2} = \frac{y^2 - x^2}{(x^2+y^2)^2}.
\]
\[
f_{yx} = \frac{\partial}{\partial x}\left(\frac{x}{x^2+y^2}\right) = \frac{(x^2+y^2) - x \cdot 2x}{(x^2+y^2)^2} = \frac{y^2 - x^2}{(x^2+y^2)^2}.
\]

Note: $f_{xx} + f_{yy} = 0$, confirming $f$ is harmonic.

% Problem 25
\item \textbf{Solution.}
$f(x,y) = \dfrac{xy}{x^2+y^2}$ for $(x,y) \neq (0,0)$ and $f(0,0) = 0$.

\textbf{Partial derivatives at $(0,0)$:}
\[
f_x(0,0) = \lim_{h \to 0}\frac{f(h,0)-f(0,0)}{h} = \lim_{h \to 0}\frac{0 - 0}{h} = 0.
\]
\[
f_y(0,0) = \lim_{k \to 0}\frac{f(0,k)-f(0,0)}{k} = \lim_{k \to 0}\frac{0 - 0}{k} = 0.
\]
So both partial derivatives exist at $(0,0)$ and equal 0.

\textbf{Not differentiable at $(0,0)$:}
If $f$ were differentiable at $(0,0)$, then
\[
\lim_{(h,k)\to(0,0)} \frac{f(h,k) - f(0,0) - f_x(0,0)h - f_y(0,0)k}{\sqrt{h^2+k^2}} = 0.
\]
This becomes:
\[
\lim_{(h,k)\to(0,0)} \frac{hk}{(h^2+k^2)\sqrt{h^2+k^2}} = \lim_{(h,k)\to(0,0)} \frac{hk}{(h^2+k^2)^{3/2}}.
\]
Along $h = k$: $\dfrac{h^2}{(2h^2)^{3/2}} = \dfrac{1}{2\sqrt{2}|h|} \to \infty$.

Since this limit does not equal 0, $f$ is not differentiable at $(0,0)$. \qed

% Problem 26
\item \textbf{Solution.}
$f(x,y) = x^3y + xy^3$. Note $f$ is homogeneous of degree 4 (each term has total degree 4).
\[
f_x = 3x^2y + y^3, \qquad f_y = x^3 + 3xy^2.
\]
\begin{align*}
xf_x + yf_y &= x(3x^2y + y^3) + y(x^3 + 3xy^2) \\
&= 3x^3y + xy^3 + x^3y + 3xy^3 \\
&= 4x^3y + 4xy^3 = 4(x^3y + xy^3) = 4f.
\end{align*}
Euler's theorem is verified. \qed

% Problem 27
\item \textbf{Solution.}
Given $f_{xx} + f_{yy} = 0$ (Laplace's equation), let $u = f_x$ and $v = f_y$.

\textbf{First identity:} $\dfrac{\partial u}{\partial x} + \dfrac{\partial v}{\partial y} = f_{xx} + f_{yy} = 0$.

\textbf{Second identity:} By Clairaut's theorem (since $f$ is twice differentiable):
\[
\frac{\partial u}{\partial y} = f_{xy} = f_{yx} = \frac{\partial v}{\partial x}.
\]
\qed

(Note: The first identity is the divergence-free condition, and the second is a consequence of equality of mixed partials --- these together resemble the Cauchy--Riemann equations.)

% Problem 28
\item \textbf{Solution.}
Given $xf_x + yf_y = nf$, define $g(t) = f(tx, ty)$ for fixed $x,y$ and $t > 0$.

By the chain rule:
\[
g'(t) = xf_x(tx,ty) + yf_y(tx,ty).
\]
Setting $u = tx$, $v = ty$, the hypothesis $uf_u + vf_v = nf$ gives:
\[
(tx)f_x(tx,ty) + (ty)f_y(tx,ty) = nf(tx,ty) \implies t\,g'(t) = n\,g(t).
\]
This is the ODE $g'(t) = \dfrac{n}{t}g(t)$, which separates to:
\[
\frac{g'}{g} = \frac{n}{t} \implies \ln|g(t)| = n\ln t + C \implies g(t) = At^n
\]
for some constant $A$. Setting $t = 1$: $g(1) = f(x,y) = A$. Therefore:
\[
g(t) = f(tx,ty) = t^n f(x,y). \quad \qed
\]

% Problem 29
\item \textbf{Solution.}
The equation is $F(x,y) = y - x^2 - \sin(xy) = 0$.

At $(0,0)$: $0 - 0 - \sin(0) = 0$.

\textbf{First derivative:}
\[
F_x = -2x - y\cos(xy), \qquad F_y = 1 - x\cos(xy).
\]
\[
\frac{dy}{dx} = -\frac{F_x}{F_y} = \frac{2x + y\cos(xy)}{1 - x\cos(xy)}.
\]
At $(0,0)$: $\dfrac{dy}{dx}\Big|_{(0,0)} = \dfrac{0 + 0}{1 - 0} = \boxed{0}$.

\textbf{Second derivative:} Let $y' = \dfrac{2x + y\cos(xy)}{1 - x\cos(xy)}$. Differentiating the original equation $y - x^2 = \sin(xy)$ implicitly again:

$y' - 2x = \cos(xy)(y + xy')$. At $(0,0)$ with $y' = 0$:
$0 - 0 = \cos(0)(0 + 0) = 0$.

Differentiating once more: $y'' - 2 = -\sin(xy)(y+xy')^2 + \cos(xy)(y' + y' + xy'') $.

At $(0,0)$: $y'' - 2 = 0 + \cos(0)(0 + 0 + 0) = 0$.

Therefore $\dfrac{d^2y}{dx^2}\Big|_{(0,0)} = \boxed{2}$.

% Problem 30
\item \textbf{Solution.}
Let $F(x,y,z) = e^{xyz} - xyz - 1 = 0$.

At any point where $xyz = 0$: $e^0 - 0 - 1 = 0$.

By implicit differentiation, treating $z$ as a function of $x$ and $y$:
\[
F_x = yze^{xyz} - yz, \qquad F_y = xze^{xyz} - xz, \qquad F_z = xye^{xyz} - xy.
\]
At $xyz = 0$: $F_x = yz(1) - yz = 0$, $F_y = xz - xz = 0$, $F_z = xy - xy = 0$.

All partials of $F$ vanish, so we cannot directly use $\partial z/\partial x = -F_x/F_z$. We use the cyclic relation for implicit functions instead.

Consider the constraint $F(x,y,z)=0$ defining a surface. For any implicit function on a smooth surface, the cyclic relation holds:
\[
\frac{\partial z}{\partial x}\bigg|_y \cdot \frac{\partial x}{\partial y}\bigg|_z \cdot \frac{\partial y}{\partial z}\bigg|_x = -1,
\]
provided the partial derivatives exist and are well-defined. This is a general identity:
\[
\frac{\partial z}{\partial x} = -\frac{F_x}{F_z}, \quad \frac{\partial x}{\partial y} = -\frac{F_y}{F_x}, \quad \frac{\partial y}{\partial z} = -\frac{F_z}{F_y}.
\]
Multiplying:
\[
\left(-\frac{F_x}{F_z}\right)\left(-\frac{F_y}{F_x}\right)\left(-\frac{F_z}{F_y}\right) = -\frac{F_x F_y F_z}{F_z F_x F_y} = -1.
\]

At points where $xyz = 0$ but the individual partials are not all zero simultaneously (e.g., at a point where exactly one of $x,y,z$ is zero but the others are nonzero), the formula applies directly.

For the case where two or more of $x,y,z$ are zero (making all $F_x = F_y = F_z = 0$), the result still holds by the general cyclic identity for smooth implicit surfaces, provided the surface is non-degenerate.

Therefore $\dfrac{\partial z}{\partial x}\cdot\dfrac{\partial x}{\partial y}\cdot\dfrac{\partial y}{\partial z} = -1$. \qed

%%% THE CHAIN RULE (Problems 31--39) %%%

% Problem 31
\item \textbf{Solution.}
$z = x^2 + y^2$, $x = t\cos t$, $y = t\sin t$.

Note $x^2 + y^2 = t^2\cos^2 t + t^2\sin^2 t = t^2$, so $z = t^2$.

Therefore $\dfrac{dz}{dt} = \boxed{2t}$.

\textbf{Verification via Chain Rule:}
\[
\frac{dz}{dt} = 2x\frac{dx}{dt} + 2y\frac{dy}{dt}.
\]
$\dfrac{dx}{dt} = \cos t - t\sin t$, $\dfrac{dy}{dt} = \sin t + t\cos t$.
\begin{align*}
\frac{dz}{dt} &= 2t\cos t(\cos t - t\sin t) + 2t\sin t(\sin t + t\cos t)\\
&= 2t(\cos^2 t - t\sin t\cos t + \sin^2 t + t\sin t\cos t) = 2t(1) = 2t.
\end{align*}

% Problem 32
\item \textbf{Solution.}
$F(x,y) = x^2y + y^3 - 8 = 0$.
\[
F_x = 2xy, \qquad F_y = x^2 + 3y^2.
\]
\[
\frac{dy}{dx} = -\frac{F_x}{F_y} = \boxed{-\frac{2xy}{x^2 + 3y^2}}.
\]

% Problem 33
\item \textbf{Solution.}
$z = xy$, $x = s^2t$, $y = st^2$.

\begin{align*}
\frac{\partial z}{\partial s} &= \frac{\partial z}{\partial x}\frac{\partial x}{\partial s} + \frac{\partial z}{\partial y}\frac{\partial y}{\partial s} = y(2st) + x(t^2) \\
&= (st^2)(2st) + (s^2t)(t^2) = 2s^2t^3 + s^2t^3 = 3s^2t^3.
\end{align*}

\begin{align*}
\frac{\partial z}{\partial t} &= \frac{\partial z}{\partial x}\frac{\partial x}{\partial t} + \frac{\partial z}{\partial y}\frac{\partial y}{\partial t} = y(s^2) + x(2st) \\
&= (st^2)(s^2) + (s^2t)(2st) = s^3t^2 + 2s^3t^2 = 3s^3t^2.
\end{align*}

\textbf{Verification:} $z = xy = (s^2t)(st^2) = s^3t^3$. Then $z_s = 3s^2t^3$ and $z_t = 3s^3t^2$, confirming the result.

% Problem 34
\item \textbf{Solution.}
$F(x,y,z) = x^2 + y^2 + z^2 - 16 = 0$.
\[
F_x = 2x, \quad F_y = 2y, \quad F_z = 2z.
\]
\[
\frac{\partial z}{\partial x} = -\frac{F_x}{F_z} = -\frac{2x}{2z} = \boxed{-\frac{x}{z}}, \qquad \frac{\partial z}{\partial y} = -\frac{F_y}{F_z} = -\frac{2y}{2z} = \boxed{-\frac{y}{z}}.
\]
(Valid when $z \neq 0$.)

% Problem 35
\item \textbf{Solution.}
$z = f(u,v)$, $u = xy$, $v = y/x$.

By the chain rule:
\[
\frac{\partial z}{\partial x} = \frac{\partial f}{\partial u}\frac{\partial u}{\partial x} + \frac{\partial f}{\partial v}\frac{\partial v}{\partial x} = f_u \cdot y + f_v \cdot \left(-\frac{y}{x^2}\right) = \boxed{yf_u - \frac{y}{x^2}f_v}.
\]

% Problem 36
\item \textbf{Solution.}
Let $I(x) = \displaystyle\int_0^x te^{xt}\,dt$. By the Leibniz rule:
\[
\frac{dI}{dx} = \underbrace{x e^{x \cdot x}}_{= x e^{x^2}} + \int_0^x \frac{\partial}{\partial x}(te^{xt})\,dt = xe^{x^2} + \int_0^x t^2 e^{xt}\,dt.
\]

Evaluate $\displaystyle\int_0^x t^2 e^{xt}\,dt$ by integration by parts (with $u = t^2$, $dv = e^{xt}\,dt$):

Let us compute $\int_0^x t^2 e^{xt}\,dt$. Substituting or using repeated integration by parts:

$\displaystyle\int t^2 e^{xt}\,dt = \frac{t^2 e^{xt}}{x} - \frac{2t e^{xt}}{x^2} + \frac{2e^{xt}}{x^3} + C$.

Evaluating from 0 to $x$:
\[
\left[\frac{t^2 e^{xt}}{x} - \frac{2te^{xt}}{x^2} + \frac{2e^{xt}}{x^3}\right]_0^x = \frac{x^2 e^{x^2}}{x} - \frac{2xe^{x^2}}{x^2} + \frac{2e^{x^2}}{x^3} - \frac{2}{x^3}
\]
\[
= xe^{x^2} - \frac{2e^{x^2}}{x} + \frac{2e^{x^2}}{x^3} - \frac{2}{x^3}.
\]

Therefore:
\[
\frac{dI}{dx} = xe^{x^2} + xe^{x^2} - \frac{2e^{x^2}}{x} + \frac{2e^{x^2}}{x^3} - \frac{2}{x^3} = \boxed{2xe^{x^2} - \frac{2e^{x^2}}{x} + \frac{2(e^{x^2}-1)}{x^3}}.
\]

% Problem 37
\item \textbf{Solution.}

\begin{enumerate}
    \item $\mathbf{r}(t) = \langle t^2, \sin t, e^t\rangle$, so
    \[
    \mathbf{r}'(t) = \langle 2t, \cos t, e^t\rangle.
    \]

    \item At $t = \pi$:
    \[
    \mathbf{r}'(\pi) = \langle 2\pi, \cos\pi, e^{\pi}\rangle = \langle 2\pi, -1, e^{\pi}\rangle.
    \]
    \begin{align*}
    |\mathbf{r}'(\pi)| &= \sqrt{(2\pi)^2 + (-1)^2 + (e^{\pi})^2}
    = \sqrt{4\pi^2 + 1 + e^{2\pi}} \\
    &\approx \sqrt{39.478 + 1 + 535.49}
    \approx \sqrt{575.97} \approx \boxed{24.0}.
    \end{align*}
\end{enumerate}

% Problem 38
\item \textbf{Solution.}
$f(x,y) = \ln(xy) + e^{x^2 y}$.

\begin{enumerate}
    \item $f_x = \dfrac{1}{x} + 2xye^{x^2y}$, \quad $f_y = \dfrac{1}{y} + x^2 e^{x^2y}$.

    At $(1,1)$: $f_x(1,1) = 1 + 2e = 1 + 2e$, \quad $f_y(1,1) = 1 + e$.

    The total differential is:
    \[
    df = f_x\,dx + f_y\,dy = (1 + 2e)\,dx + (1 + e)\,dy.
    \]

    \item With $dx = 0.01$, $dy = -0.01$:
    \[
    \Delta f \approx (1 + 2e)(0.01) + (1+e)(-0.01) = 0.01[(1+2e) - (1+e)] = 0.01 \cdot e \approx 0.0272.
    \]
\end{enumerate}

% Problem 39
\item \textbf{Solution.}
$f(x,y) = \displaystyle\sum_{k=1}^{3}\sin(x^k y^{3-k}) = \sin(xy^2) + \sin(x^2y) + \sin(x^3)$.

By the chain rule:
\[
\frac{df}{dt} = f_x \frac{dx}{dt} + f_y \frac{dy}{dt}.
\]

Compute the partials:
\[
f_x = y^2\cos(xy^2) + 2xy\cos(x^2y) + 3x^2\cos(x^3).
\]
\[
f_y = 2xy\cos(xy^2) + x^2\cos(x^2y).
\]

Therefore:
\[
\frac{df}{dt} = \bigl[y^2\cos(xy^2) + 2xy\cos(x^2y) + 3x^2\cos(x^3)\bigr]x'(t) + \bigl[2xy\cos(xy^2) + x^2\cos(x^2y)\bigr]y'(t).
\]

More compactly, using the general term:
\[
\frac{df}{dt} = \sum_{k=1}^{3}\cos(x^k y^{3-k})\left[kx^{k-1}y^{3-k}x'(t) + (3-k)x^k y^{2-k}y'(t)\right].
\]

%%% DIRECTIONAL DERIVATIVES AND GRADIENT VECTORS (Problems 40--47) %%%

% Problem 40
\item \textbf{Solution.}
$f(x,y) = x^2y + y^3$.
\[
\nabla f = \langle 2xy + 0,\; x^2 + 3y^2 \rangle = \langle 2xy,\; x^2 + 3y^2\rangle.
\]
At $(1,2)$:
\[
\nabla f(1,2) = \langle 2(1)(2),\; 1 + 3(4)\rangle = \boxed{\langle 4, 13\rangle}.
\]

% Problem 41
\item \textbf{Solution.}
$f(x,y) = xe^y$, $\mathbf{v} = \langle 3, 4\rangle$.

$\nabla f = \langle e^y, xe^y\rangle$. At $(1,0)$: $\nabla f(1,0) = \langle 1, 1\rangle$.

Unit vector: $\mathbf{u} = \dfrac{\mathbf{v}}{|\mathbf{v}|} = \dfrac{1}{5}\langle 3,4\rangle$.

\[
D_{\mathbf{u}}f = \nabla f \cdot \mathbf{u} = \langle 1,1\rangle \cdot \frac{1}{5}\langle 3,4\rangle = \frac{3+4}{5} = \boxed{\frac{7}{5}}.
\]

% Problem 42
\item \textbf{Solution.}
$f(x,y) = 2x^2 + 3y^2$. The direction of steepest increase is the gradient direction.
\[
\nabla f = \langle 4x, 6y\rangle.
\]
At $(1,1)$: $\nabla f(1,1) = \langle 4, 6\rangle$.

The direction of steepest increase is $\boxed{\langle 4, 6\rangle}$, or as a unit vector: $\dfrac{1}{\sqrt{52}}\langle 4,6\rangle = \dfrac{1}{2\sqrt{13}}\langle 4,6\rangle = \left\langle \dfrac{2}{\sqrt{13}}, \dfrac{3}{\sqrt{13}}\right\rangle$.

% Problem 43
\item \textbf{Solution.}
$f(x,y,z) = xyz$, $\mathbf{u} = \langle 1,1,1\rangle$ (must normalize).
\[
\nabla f = \langle yz, xz, xy\rangle.
\]
At $(1,2,3)$: $\nabla f = \langle 6, 3, 2\rangle$.

Unit vector: $\hat{\mathbf{u}} = \dfrac{1}{\sqrt{3}}\langle 1,1,1\rangle$.

\[
D_{\hat{\mathbf{u}}}f = \langle 6,3,2\rangle \cdot \frac{1}{\sqrt{3}}\langle 1,1,1\rangle = \frac{6+3+2}{\sqrt{3}} = \boxed{\frac{11}{\sqrt{3}} = \frac{11\sqrt{3}}{3}}.
\]

% Problem 44
\item \textbf{Solution.}
$f(x,y) = \sin(xy)$.
\[
\nabla f(x,y) = \langle y\cos(xy),\; x\cos(xy)\rangle.
\]
At $(\pi, 0)$:
\[
\nabla f(\pi, 0) = \langle 0 \cdot \cos(0),\; \pi\cos(0)\rangle = \boxed{\langle 0, \pi\rangle}.
\]

% Problem 45
\item \textbf{Solution.}
$f(x,y) = \ln(x^2+y^2)$.
\[
\nabla f = \left\langle \frac{2x}{x^2+y^2},\; \frac{2y}{x^2+y^2}\right\rangle.
\]
At $(1,1)$: $\nabla f(1,1) = \left\langle \frac{2}{2}, \frac{2}{2}\right\rangle = \langle 1, 1\rangle$.

Direction from $(1,1)$ to $(2,3)$: $\mathbf{v} = \langle 1, 2\rangle$, $|\mathbf{v}| = \sqrt{5}$.

$\hat{\mathbf{u}} = \dfrac{1}{\sqrt{5}}\langle 1, 2\rangle$.

\[
D_{\hat{\mathbf{u}}}f = \langle 1,1\rangle \cdot \frac{1}{\sqrt{5}}\langle 1,2\rangle = \frac{1+2}{\sqrt{5}} = \boxed{\frac{3}{\sqrt{5}} = \frac{3\sqrt{5}}{5}}.
\]

% Problem 46
\item \textbf{Solution.}
$T(x,y) = e^{-x^2-y^2}$.
\[
\nabla T = \langle -2xe^{-x^2-y^2},\; -2ye^{-x^2-y^2}\rangle = -2e^{-x^2-y^2}\langle x, y\rangle.
\]
At $(1,1)$:
\[
\nabla T(1,1) = -2e^{-2}\langle 1,1\rangle.
\]
The direction of steepest increase is $\nabla T(1,1) = -2e^{-2}\langle 1,1\rangle$, which points in the direction $\langle -1, -1\rangle$.

As a unit vector, the direction of steepest temperature increase is $\boxed{\left\langle -\dfrac{1}{\sqrt{2}}, -\dfrac{1}{\sqrt{2}}\right\rangle}$, i.e., toward the origin. This makes physical sense: the temperature $T = e^{-x^2-y^2}$ is maximized at the origin, so at $(1,1)$ the temperature increases most rapidly toward the origin.

% Problem 47
\item \textbf{Solution.}
$J(w,b) = \dfrac{1}{2}(wx + b - y)^2$.

Let $\hat{y} = wx + b$. Then $J = \frac{1}{2}(\hat{y} - y)^2$.

\[
\frac{\partial J}{\partial w} = (wx + b - y) \cdot x = \boxed{x(wx + b - y)}.
\]
\[
\frac{\partial J}{\partial b} = (wx + b - y) \cdot 1 = \boxed{wx + b - y}.
\]

These are the gradients used in gradient descent for linear regression.

%%% TANGENT PLANES AND LINEAR APPROXIMATIONS (Problems 48--53) %%%

% Problem 48
\item \textbf{Solution.}
$z = f(x,y) = x^2 + y^2$ at $(1,1,2)$.
\[
f_x = 2x, \quad f_y = 2y.
\]
At $(1,1)$: $f_x = 2$, $f_y = 2$.

Tangent plane:
\[
z - 2 = 2(x-1) + 2(y-1) \implies \boxed{z = 2x + 2y - 2}.
\]

% Problem 49
\item \textbf{Solution.}
We want to approximate $\sqrt{4.1^2 + 2.9^2}$. Let $f(x,y) = \sqrt{x^2+y^2}$ near $(4, 3)$ (where $f(4,3) = 5$).
\[
f_x = \frac{x}{\sqrt{x^2+y^2}}, \quad f_y = \frac{y}{\sqrt{x^2+y^2}}.
\]
At $(4,3)$: $f_x = \frac{4}{5}$, $f_y = \frac{3}{5}$.

With $\Delta x = 0.1$, $\Delta y = -0.1$:
\[
f(4.1, 2.9) \approx 5 + \frac{4}{5}(0.1) + \frac{3}{5}(-0.1) = 5 + 0.08 - 0.06 = \boxed{5.02}.
\]

% Problem 50
\item \textbf{Solution.}
$f(x,y) = e^x\sin y$ near $(0,0)$. Note $f(0,0) = 0$.
\[
f_x = e^x\sin y, \quad f_y = e^x\cos y.
\]
At $(0,0)$: $f_x = 0$, $f_y = 1$.

Linear approximation:
\[
L(x,y) = f(0,0) + f_x(0,0)\,x + f_y(0,0)\,y = 0 + 0 \cdot x + 1 \cdot y = \boxed{y}.
\]

% Problem 51
\item \textbf{Solution.}
$z = \ln(xy)$ at $(1,1,0)$.
\[
f_x = \frac{1}{x}, \quad f_y = \frac{1}{y}.
\]
At $(1,1)$: $f_x = 1$, $f_y = 1$.

Tangent plane:
\[
z - 0 = 1(x - 1) + 1(y - 1) \implies \boxed{z = x + y - 2}.
\]

% Problem 52
\item \textbf{Solution.}
$f(x,y) = \sqrt{x^2+y^2}$ near $(1,1)$. $f(1,1) = \sqrt{2}$.
\[
f_x = \frac{x}{\sqrt{x^2+y^2}}, \quad f_y = \frac{y}{\sqrt{x^2+y^2}}.
\]
At $(1,1)$: $f_x = f_y = \dfrac{1}{\sqrt{2}}$.

\[
f(1.1, 0.9) \approx \sqrt{2} + \frac{1}{\sqrt{2}}(0.1) + \frac{1}{\sqrt{2}}(-0.1) = \sqrt{2} + 0 = \boxed{\sqrt{2} \approx 1.414}.
\]

% Problem 53
\item \textbf{Solution.}
$z = xe^y$ at $(1,0,1)$.
\[
f_x = e^y, \quad f_y = xe^y.
\]
At $(1,0)$: $f_x = 1$, $f_y = 1$.

Tangent plane:
\[
z - 1 = 1(x-1) + 1(y - 0) \implies \boxed{z = x + y}.
\]

\textbf{Above or below?} Consider $g(x,y) = xe^y - (x+y)$, the difference between the surface and the plane.

$g(1,0) = 0$. The Hessian of $f$ at $(1,0)$:
\[
f_{xx} = 0, \quad f_{xy} = e^y\big|_{(1,0)} = 1, \quad f_{yy} = xe^y\big|_{(1,0)} = 1.
\]
The Hessian matrix is $H = \begin{pmatrix} 0 & 1 \\ 1 & 1\end{pmatrix}$ with determinant $-1 < 0$.

Since the Hessian is indefinite, the surface is neither entirely above nor entirely below the tangent plane near $(1,0,1)$ --- it lies above in some directions and below in others (a saddle-type behavior).

%%% EXTREMA AND CRITICAL POINTS (Problems 54--61) %%%

% Problem 54
\item \textbf{Solution.}
$f(x,y) = x^2 + y^2 - 4x - 6y$.
\[
f_x = 2x - 4 = 0 \implies x = 2.
\]
\[
f_y = 2y - 6 = 0 \implies y = 3.
\]
The only critical point is $\boxed{(2, 3)}$.

Classification: $f_{xx} = 2$, $f_{yy} = 2$, $f_{xy} = 0$. $D = f_{xx}f_{yy} - f_{xy}^2 = 4 > 0$ and $f_{xx} > 0$, so $(2,3)$ is a \textbf{local minimum} with $f(2,3) = 4 + 9 - 8 - 18 = -13$.

% Problem 55
\item \textbf{Solution.}
$f(x,y) = x^2y - \dfrac{y^3}{3}$.
\[
f_x = 2xy = 0, \qquad f_y = x^2 - y^2 = 0.
\]
From $f_x = 0$: $x = 0$ or $y = 0$.

\textbf{Case 1:} $x = 0$. Then $f_y = -y^2 = 0 \implies y = 0$. Critical point: $(0,0)$.

\textbf{Case 2:} $y = 0$. Then $f_y = x^2 = 0 \implies x = 0$. Same point $(0,0)$.

So the only critical point is $(0,0)$.

Second derivatives: $f_{xx} = 2y$, $f_{yy} = -2y$, $f_{xy} = 2x$.

At $(0,0)$: $f_{xx} = 0$, $f_{yy} = 0$, $f_{xy} = 0$. $D = 0$: \textbf{inconclusive}.

Examine $f$ directly: $f(x,0) = 0$ for all $x$, $f(0,y) = -y^3/3$ (changes sign), and $f(\epsilon,\epsilon) = \epsilon^3 - \epsilon^3/3 = 2\epsilon^3/3$ (changes sign). Since $f$ takes both positive and negative values near $(0,0)$, and $f(0,0) = 0$, the origin is a \textbf{saddle point}.

% Problem 56
\item \textbf{Solution.}
$f(x,y) = x^2 + y^2$ on $x^2 + y^2 \leq 4$.

\textbf{Interior:} $f_x = 2x = 0$, $f_y = 2y = 0 \implies (0,0)$. $f(0,0) = 0$.

\textbf{Boundary:} $x^2 + y^2 = 4$, so $f = 4$ everywhere on the boundary.

\textbf{Absolute minimum:} $f(0,0) = \boxed{0}$.

\textbf{Absolute maximum:} $f = \boxed{4}$ on the circle $x^2 + y^2 = 4$.

% Problem 57
\item \textbf{Solution.}
$f(x,y) = xe^y - x^2 - y^2$.
\[
f_x = e^y - 2x = 0 \implies x = \frac{e^y}{2}.
\]
\[
f_y = xe^y - 2y = 0 \implies \frac{e^y}{2}\cdot e^y = 2y \implies e^{2y} = 4y.
\]
Define $h(y) = e^{2y} - 4y$. Then $h(0) = 1 > 0$ and $h'(y) = 2e^{2y} - 4$. Setting $h'(y) = 0$ gives $e^{2y} = 2$, i.e., $y = \frac{\ln 2}{2} \approx 0.347$. The minimum value is $h(\tfrac{\ln 2}{2}) = 2 - 2\ln 2 \approx 0.614 > 0$. For $y < 0$, $4y < 0 < e^{2y}$, so no solution there either. Since $h(y) > 0$ for all $y$, the equation $e^{2y} = 4y$ has \textbf{no real solution}, and therefore there are \textbf{no critical points}.

Moreover, $f$ is unbounded above: along the curve $x = e^y/2$, we have $f = \frac{e^{2y}}{2} - \frac{e^{2y}}{4} - y^2 = \frac{e^{2y}}{4} - y^2 \to \infty$ as $y \to \infty$. Similarly, $f$ is unbounded below (along $x \to -\infty$ with fixed $y$). Therefore $f$ has \textbf{no critical points} and \textbf{no local extrema}.

% Problem 58
\item \textbf{Solution.}
$f(x,y) = xy - x^2 - y^2$.
\[
f_x = y - 2x = 0, \qquad f_y = x - 2y = 0.
\]
From the first: $y = 2x$. Substituting: $x - 4x = 0 \implies x = 0$, hence $y = 0$.

Critical point: $(0,0)$.

$f_{xx} = -2$, $f_{yy} = -2$, $f_{xy} = 1$. $D = (-2)(-2) - 1^2 = 3 > 0$ and $f_{xx} < 0$.

$(0,0)$ is a \textbf{local maximum} with $f(0,0) = \boxed{0}$.

% Problem 59
\item \textbf{Solution.}
$f(x,y) = (y - x^2)(y - 2x^2) = y^2 - 3x^2y + 2x^4$.
\[
f_x = -6xy + 8x^3 = 2x(4x^2 - 3y) = 0.
\]
\[
f_y = 2y - 3x^2 = 0.
\]
From $f_x = 0$: $x = 0$ or $y = \frac{4x^2}{3}$.

\textbf{Case 1:} $x = 0$. Then $f_y = 2y = 0 \implies y = 0$. Critical point: $(0,0)$.

\textbf{Case 2:} $y = \frac{4x^2}{3}$. Then $f_y = \frac{8x^2}{3} - 3x^2 = -\frac{x^2}{3} = 0 \implies x = 0$, giving $(0,0)$ again.

So the \textbf{only critical point} is $(0,0)$.

\textbf{Second derivative test:}
$f_{xx} = -6y + 24x^2$, $f_{yy} = 2$, $f_{xy} = -6x$.
At $(0,0)$: $f_{xx} = 0$, $f_{yy} = 2$, $f_{xy} = 0$. $D = 0$: \textbf{inconclusive}.

\textbf{Direct analysis:}
$f(0,0) = 0$. Along $y = cx^2$ for constant $c$:
\[
f = (cx^2 - x^2)(cx^2 - 2x^2) = x^4(c-1)(c-2).
\]
\begin{itemize}
    \item If $c < 1$ or $c > 2$: $(c-1)(c-2) > 0$, so $f > 0$ near origin.
    \item If $1 < c < 2$: $(c-1)(c-2) < 0$, so $f < 0$ near origin.
\end{itemize}

Since $f$ takes both positive and negative values in every neighborhood of $(0,0)$, and $f(0,0) = 0$, the origin is a \textbf{saddle point}.

% Problem 60
\item \textbf{Solution.}
$f(x,y,z) = x^2 + y^2 + z^2 - 4x - 6y + 8z$.

\begin{enumerate}
    \item Setting all partial derivatives to zero:
    \[
    f_x = 2x - 4 = 0 \implies x = 2.
    \]
    \[
    f_y = 2y - 6 = 0 \implies y = 3.
    \]
    \[
    f_z = 2z + 8 = 0 \implies z = -4.
    \]
    The only critical point is $\boxed{(2, 3, -4)}$.

    \item The Hessian matrix is:
    \[
    H = \begin{pmatrix} 2 & 0 & 0 \\ 0 & 2 & 0 \\ 0 & 0 & 2 \end{pmatrix}.
    \]
    All eigenvalues are 2 (positive), so $H$ is positive definite. Therefore $(2,3,-4)$ is a \textbf{local minimum} with $f(2,3,-4) = 4 + 9 + 16 - 8 - 18 - 32 = -29$.
\end{enumerate}

% Problem 61
\item \textbf{Solution.}
For $f(\mathbf{x})$ with $\mathbf{x} \in \mathbb{R}^n$, a critical point $\mathbf{x}_0$ satisfies $\nabla f(\mathbf{x}_0) = \mathbf{0}$.

The classification depends on the \textbf{Hessian matrix} $H(\mathbf{x}_0) = \left[\dfrac{\partial^2 f}{\partial x_i \partial x_j}\right]_{n \times n}$:

\begin{itemize}
    \item \textbf{Local minimum:} $H(\mathbf{x}_0)$ is \textbf{positive definite} (all eigenvalues $> 0$).
    \item \textbf{Local maximum:} $H(\mathbf{x}_0)$ is \textbf{negative definite} (all eigenvalues $< 0$).
    \item \textbf{Saddle point:} $H(\mathbf{x}_0)$ is \textbf{indefinite} (has both positive and negative eigenvalues).
    \item \textbf{Inconclusive:} $H(\mathbf{x}_0)$ is singular (has a zero eigenvalue).
\end{itemize}

\textbf{Extension from two variables:} In the two-variable case, $H = \begin{pmatrix} f_{xx} & f_{xy} \\ f_{xy} & f_{yy}\end{pmatrix}$, and positive definiteness is equivalent to $D = f_{xx}f_{yy} - f_{xy}^2 > 0$ with $f_{xx} > 0$ (by Sylvester's criterion: all leading principal minors are positive). For $n$ variables, Sylvester's criterion generalizes: $H$ is positive definite if and only if all $n$ leading principal minors are positive, i.e.,
\[
H_1 > 0, \quad H_2 > 0, \quad \ldots, \quad H_n > 0,
\]
where $H_k = \det\!\left[\frac{\partial^2 f}{\partial x_i \partial x_j}\right]_{1 \leq i,j \leq k}$.

Alternatively (and often more practically), one computes the eigenvalues of $H$.

%%% OPTIMIZATION (Problems 62--71) %%%

% Problem 62
\item \textbf{Solution.}
Distance from $(2,1,3)$ to the plane $x + y + z = 1$.

Using the distance formula:
\[
d = \frac{|ax_0 + by_0 + cz_0 - d|}{\sqrt{a^2+b^2+c^2}} = \frac{|2 + 1 + 3 - 1|}{\sqrt{1+1+1}} = \frac{5}{\sqrt{3}} = \boxed{\frac{5\sqrt{3}}{3}}.
\]

% Problem 63
\item \textbf{Solution.}
Minimize $S = 2(xy + yz + xz)$ subject to $V = xyz = 100$.

By Lagrange multipliers (or by symmetry): a box of minimum surface area with fixed volume is a \textbf{cube}. So $x = y = z = \sqrt[3]{100} = 100^{1/3}$.

\textbf{Verification with Lagrange multipliers:} $\nabla S = \lambda \nabla V$:
\[
2(y+z) = \lambda yz, \quad 2(x+z) = \lambda xz, \quad 2(x+y) = \lambda xy.
\]
From the first two: $\dfrac{2(y+z)}{yz} = \dfrac{2(x+z)}{xz}$, giving $\dfrac{1}{z}+\dfrac{1}{y} = \dfrac{1}{z}+\dfrac{1}{x}$, so $x = y$. Similarly $y = z$.

With $x = y = z$: $x^3 = 100$, so $x = y = z = \boxed{\sqrt[3]{100} \approx 4.642}$ units.

The minimum surface area is $S = 6(\sqrt[3]{100})^2 = 6 \cdot 100^{2/3} \approx 129.3$ square units.

% Problem 64
\item \textbf{Solution.}
Minimize $D^2 = x^2 + (x^2 - 1)^2$ (distance squared from $(x, x^2)$ to $(0,1)$).

Let $g(x) = x^2 + (x^2-1)^2 = x^2 + x^4 - 2x^2 + 1 = x^4 - x^2 + 1$.
\[
g'(x) = 4x^3 - 2x = 2x(2x^2-1) = 0.
\]
Solutions: $x = 0$ or $x = \pm\dfrac{1}{\sqrt{2}}$.

\[
g(0) = 1, \qquad g\!\left(\pm\frac{1}{\sqrt{2}}\right) = \frac{1}{4} - \frac{1}{2} + 1 = \frac{3}{4}.
\]

The minimum distance squared is $\frac{3}{4}$, so the minimum distance is $\dfrac{\sqrt{3}}{2}$.

The closest points are $\boxed{\left(\pm\dfrac{1}{\sqrt{2}},\; \dfrac{1}{2}\right)}$.

% Problem 65
\item \textbf{Solution.}
Minimize $f(x,y) = x^2 + y^2$ subject to $g(x,y) = x^3 + y^3 - 8 = 0$.

Lagrange conditions: $\nabla f = \lambda \nabla g$:
\[
2x = 3\lambda x^2, \qquad 2y = 3\lambda y^2.
\]

\textbf{Case 1:} $x \neq 0$ and $y \neq 0$. Then $\lambda = \dfrac{2}{3x}$ and $\lambda = \dfrac{2}{3y}$, giving $x = y$. Substituting into the constraint: $2x^3 = 8 \implies x^3 = 4 \implies x = y = \sqrt[3]{4}$.

$f = 2(\sqrt[3]{4})^2 = 2 \cdot 4^{2/3} = 2\sqrt[3]{16} \approx 5.04$.

\textbf{Case 2:} $x = 0$. Then $y^3 = 8 \implies y = 2$. $f = 4$.

\textbf{Case 3:} $y = 0$. Then $x^3 = 8 \implies x = 2$. $f = 4$.

The minimum value is $\boxed{f = 4}$, achieved at $(2,0)$ and $(0,2)$.

% Problem 66
\item \textbf{Solution.}
Optimize $f(x,y) = xy$ subject to $g(x,y) = x^2 + y^2 - 16 = 0$.

$\nabla f = \lambda \nabla g$: $y = 2\lambda x$, $x = 2\lambda y$.

Multiplying: $xy = 4\lambda^2 xy$.

\textbf{Case 1:} $xy \neq 0$. Then $4\lambda^2 = 1$, so $\lambda = \pm\frac{1}{2}$.

If $\lambda = \frac{1}{2}$: $y = x$. Constraint: $2x^2 = 16 \implies x = \pm 2\sqrt{2}$. Points: $(2\sqrt{2}, 2\sqrt{2})$ and $(-2\sqrt{2}, -2\sqrt{2})$. $f = 8$.

If $\lambda = -\frac{1}{2}$: $y = -x$. Constraint: $2x^2 = 16 \implies x = \pm 2\sqrt{2}$. Points: $(2\sqrt{2}, -2\sqrt{2})$ and $(-2\sqrt{2}, 2\sqrt{2})$. $f = -8$.

\textbf{Case 2:} $xy = 0$. Then one variable is 0, and the other is $\pm 4$. $f = 0$.

\[
\boxed{\text{Maximum} = 8, \qquad \text{Minimum} = -8.}
\]

% Problem 67
\item \textbf{Solution.}
$f(x,y,z) = x^2 + y^2 - z^2$ subject to $g(x,y,z) = x + y + z - 1 = 0$.

\begin{enumerate}
    \item $\nabla f = \lambda \nabla g$:
    \[
    2x = \lambda, \quad 2y = \lambda, \quad -2z = \lambda.
    \]
    From the first two: $x = y$. From the first and third: $2x = -(-2z) \implies z = -x$.

    Constraint: $x + x + (-x) = 1 \implies x = 1$, so $y = 1$, $z = -1$.

    Critical point: $\boxed{(1, 1, -1)}$ with $\lambda = 2$ and $f(1,1,-1) = 1 + 1 - 1 = 1$.

    \item Since we optimize on a plane (unbounded constraint set), we check whether this is a max or min. Parameterize the constraint: $z = 1 - x - y$, then
    \[
    f = x^2 + y^2 - (1-x-y)^2 = x^2 + y^2 - 1 + 2x + 2y - x^2 - 2xy - y^2 = -1 + 2x + 2y - 2xy.
    \]
    This is an indefinite quadratic in $x$ and $y$ (the $-2xy$ term makes it a saddle). More carefully, the Hessian of this reduced function is $\begin{pmatrix} 0 & -2 \\ -2 & 0 \end{pmatrix}$ with eigenvalues $\pm 2$, confirming the critical point is a \textbf{saddle point} of $f$ restricted to the constraint.

    Alternatively, note $f \to +\infty$ (take $x$ large, $z = 1-x-y$) and $f \to -\infty$ (take $z$ large, $x = y = \frac{1-z}{2}$). The constrained function has neither a maximum nor a minimum; $(1,1,-1)$ is a \textbf{saddle point}.
\end{enumerate}

% Problem 68
\item \textbf{Solution.}
Optimize $f(x,y,z) = x + y + z$ subject to $x^2 + y^2 + z^2 = 1$.

$\nabla f = \lambda \nabla g$: $1 = 2\lambda x$, $1 = 2\lambda y$, $1 = 2\lambda z$.

So $x = y = z = \dfrac{1}{2\lambda}$. Constraint: $\dfrac{3}{4\lambda^2} = 1 \implies \lambda = \pm\dfrac{\sqrt{3}}{2}$.

$\lambda = \frac{\sqrt{3}}{2}$: $x = y = z = \frac{1}{\sqrt{3}}$. $f = \frac{3}{\sqrt{3}} = \sqrt{3}$.

$\lambda = -\frac{\sqrt{3}}{2}$: $x = y = z = -\frac{1}{\sqrt{3}}$. $f = -\sqrt{3}$.

\[
\boxed{\text{Maximum} = \sqrt{3}, \qquad \text{Minimum} = -\sqrt{3}.}
\]

% Problem 69
\item \textbf{Solution.}
Minimize $f(x,y) = 3x + 4y$ subject to $x + 2y \geq 8$ and $x - y \geq 2$.

This is a linear programming problem. The minimum of a linear objective on a convex polyhedral region occurs at a vertex.

\textbf{Find vertices of the feasible region:}

Intersection of $x + 2y = 8$ and $x - y = 2$: Subtracting: $3y = 6 \implies y = 2$, $x = 4$. Point: $(4, 2)$.

Check feasibility: $4 + 4 = 8 \geq 8$ and $4 - 2 = 2 \geq 2$, so both constraints are satisfied.

The feasible region is unbounded. As $x \to \infty$ along either constraint boundary, $f \to \infty$. So the minimum is at a vertex.

$f(4, 2) = 12 + 8 = 20$.

Check boundary rays: Along $x + 2y = 8$ (away from vertex): as $y \to -\infty$, $x = 8 - 2y \to \infty$, need $x - y = 8 - 3y \geq 2$, i.e., $y \leq 2$. So $f = 3(8-2y) + 4y = 24 - 2y$. For $y \leq 2$, $f \geq 20$. Along $x - y = 2$: $x = y + 2$, need $y + 2 + 2y \geq 8$, i.e., $y \geq 2$. $f = 3(y+2) + 4y = 7y + 6$. For $y \geq 2$, $f \geq 20$.

The minimum value is $\boxed{f = 20}$ at $(4, 2)$.

% Problem 70
\item \textbf{Solution.}
Minimize $f(x,y,z) = x^2 + y^2 + z^2$ subject to
\[
g_1 = x + y + z - 6 \geq 0, \qquad g_2 = x \geq 0, \qquad g_3 = y \geq 0, \qquad g_4 = z \geq 0.
\]

\textbf{Stationarity.} $\nabla f = \mu_1 \nabla g_1 + \mu_2 \nabla g_2 + \mu_3 \nabla g_3 + \mu_4 \nabla g_4$:
\[
(2x,\, 2y,\, 2z) = \mu_1(1,1,1) + \mu_2(1,0,0) + \mu_3(0,1,0) + \mu_4(0,0,1),
\]
which gives the component equations
\[
2x = \mu_1 + \mu_2, \qquad 2y = \mu_1 + \mu_3, \qquad 2z = \mu_1 + \mu_4.
\]

\textbf{Try the natural case: $g_1$ active, $g_2, g_3, g_4$ inactive.} Set $\mu_2 = \mu_3 = \mu_4 = 0$ and require $x + y + z = 6$. The stationarity equations reduce to $2x = 2y = 2z = \mu_1$, so $x = y = z$. Combined with $x+y+z = 6$:
\[
x = y = z = 2, \qquad \mu_1 = 4.
\]

\textbf{Verify all four KKT conditions at $(2,2,2)$ with $\mu = (4,0,0,0)$.}
\begin{itemize}
\item \emph{Stationarity}: $(4,4,4) = 4(1,1,1) + 0 + 0 + 0$ \checkmark.
\item \emph{Primal feasibility}: $g_1 = 2+2+2-6 = 0 \geq 0$ \checkmark; $g_2 = g_3 = g_4 = 2 \geq 0$ \checkmark.
\item \emph{Dual feasibility}: $\mu_1 = 4 \geq 0$, $\mu_2 = \mu_3 = \mu_4 = 0 \geq 0$ \checkmark.
\item \emph{Complementary slackness}: $\mu_1 g_1 = 4 \cdot 0 = 0$ \checkmark; $\mu_i g_i = 0 \cdot 2 = 0$ for $i = 2,3,4$ \checkmark.
\end{itemize}

All four conditions hold. Since $f$ is strictly convex and the feasible region is convex, this KKT point is the unique global minimum. Only the sum constraint is active; the three non-negativity constraints are slack.
\[
\boxed{f_{\min} = 12 \text{ at } (x,y,z) = (2,2,2).}
\]
\textit{Geometric picture.} The level surfaces of $f$ are spheres centered at the origin. The minimum is the closest point of the half-space $x+y+z \geq 6$ (intersected with the first octant) to the origin --- the foot of the perpendicular from the origin to the plane $x+y+z = 6$, which is $(2,2,2)$. The non-negativity constraints are inactive precisely because that foot already lies in the first octant.

% Problem 71
\item \textbf{Solution.}
Minimize $f(x,y) = x^2 + y^2$ subject to $g_1(x,y) = x + y - 4 \geq 0$ and $g_2(x,y) = x - 2y + 1 \geq 0$.

\begin{enumerate}
    \item \textbf{KKT conditions.} For a minimization problem with constraints $g_i \geq 0$, the KKT conditions are:
    \[
    \nabla f = \mu_1 \nabla g_1 + \mu_2 \nabla g_2, \qquad \mu_i \geq 0, \qquad \mu_i g_i = 0.
    \]
    \[
    (2x, 2y) = \mu_1(1,1) + \mu_2(1,-2).
    \]
    This gives:
    \[
    2x = \mu_1 + \mu_2, \qquad 2y = \mu_1 - 2\mu_2.
    \]

    \textbf{Case 1: Both constraints active} ($g_1 = 0$ and $g_2 = 0$).

    $x + y = 4$ and $x - 2y = -1$. Solving: $3y = 5 \implies y = 5/3$, $x = 4 - 5/3 = 7/3$.

    Point: $(7/3, 5/3)$. Then $\mu_1 = 2(7/3) - \mu_2$ and $2(5/3) = \mu_1 - 2\mu_2$. From the first: $\mu_1 = 14/3 - \mu_2$. Substituting: $10/3 = 14/3 - \mu_2 - 2\mu_2 \implies 3\mu_2 = 4/3 \implies \mu_2 = 4/9$. Then $\mu_1 = 14/3 - 4/9 = 38/9$.

    Both $\mu_1 = 38/9 > 0$ and $\mu_2 = 4/9 > 0$, as required.

    $f(7/3, 5/3) = 49/9 + 25/9 = 74/9$.

    \textbf{Case 2: Only $g_1$ active} ($\mu_2 = 0$, $g_1 = 0$).

    $\mu_2 = 0$: $2x = \mu_1$, $2y = \mu_1$, so $x = y$. From $x + y = 4$: $x = y = 2$. Then $\mu_1 = 4 > 0$, as required.

    Check $g_2$: $2 - 4 + 1 = -1 < 0$. \textbf{Violated.} Not feasible.

    \textbf{Case 3: Only $g_2$ active} ($\mu_1 = 0$, $g_2 = 0$).

    $\mu_1 = 0$: $2x = \mu_2$, $2y = -2\mu_2$, so $y = -2x$ and $\mu_2 = 2x \geq 0$ requires $x \geq 0$.

    From $x - 2y = -1$: $x + 4x = -1 \implies 5x = -1 \implies x = -1/5$. But $x < 0$ means $\mu_2 < 0$. \textbf{Violates} $\mu_2 \geq 0$.

    \textbf{Case 4: No constraint active} ($\mu_1 = \mu_2 = 0$).

    Then $x = y = 0$, but $g_1(0,0) = -4 < 0$. \textbf{Not feasible.}

    \item The only candidate satisfying all KKT conditions is $\left(\dfrac{7}{3}, \dfrac{5}{3}\right)$.

    Since the objective $f$ is strictly convex and the feasible set is convex (intersection of half-planes), any KKT point is a global minimum. Therefore:

    \[
    \boxed{f_{\min} = \frac{74}{9} \approx 8.22 \text{ at } \left(\frac{7}{3}, \frac{5}{3}\right).}
    \]
\end{enumerate}

\end{enumerate}
